%!TEX root = ../main.tex
%----------------------------------------------------------------------
%  Appendix: Proofs
%  Paper:   Theory-Informed Kernel Ridge Regression for Asset Pricing
%  Authors: Amedeo Andriollo and Juan F. Imbet
%----------------------------------------------------------------------

\section{Proofs}
\label{app:proofs}

%----------------------------------------------------------------------
\subsection{Proof of Proposition~\ref{prop:representer} (Representer Theorem)}
%----------------------------------------------------------------------

\begin{proof}
Decompose any $f \in \mathcal{H}$ as $f = f_{\parallel} + f_{\perp}$,
where $f_{\parallel}$ lies in the finite-dimensional subspace
$\mathcal{V} = \operatorname{span}\{k(\mathbf{X}_s, \cdot) : s = 1, \ldots, n\}$
and $f_{\perp} \in \mathcal{V}^{\perp}$.
By the reproducing property,
\[
f(\mathbf{X}_s)
= \langle f, k(\mathbf{X}_s, \cdot) \rangle_{\mathcal{H}}
= \langle f_{\parallel}, k(\mathbf{X}_s, \cdot) \rangle_{\mathcal{H}}
  + \langle f_{\perp}, k(\mathbf{X}_s, \cdot) \rangle_{\mathcal{H}}
= f_{\parallel}(\mathbf{X}_s),
\]
since $k(\mathbf{X}_s, \cdot) \in \mathcal{V}$ and
$f_{\perp} \perp \mathcal{V}$.
Therefore the prediction loss
$\frac{1}{n}\sum_s (R_s - f(\mathbf{X}_s))^2$
depends only on $f_{\parallel}$, and so does every $C_j(f)$
by assumption~\eqref{eq:penalty_eval}.

For the RKHS norm, orthogonality gives
$\|f\|_{\mathcal{H}}^2 = \|f_{\parallel}\|_{\mathcal{H}}^2 +
\|f_{\perp}\|_{\mathcal{H}}^2$.
Since $\|f_{\perp}\|_{\mathcal{H}}^2 \geq 0$ and
$\lambda_{\mathrm{stat}} > 0$, the objective is strictly decreased
by setting $f_{\perp} = 0$ whenever $f_{\perp} \neq 0$.
Hence any minimizer satisfies $f = f_{\parallel} \in \mathcal{V}$,
which yields the representation~\eqref{eq:representer}.
\end{proof}

%----------------------------------------------------------------------
\subsection{Proof of Proposition~\ref{prop:finite_dim} (Finite-Dimensional Problem)}
%----------------------------------------------------------------------

\begin{proof}
The prediction loss becomes
$\frac{1}{n}\|\mathbf{R} - \mathbf{K}\boldsymbol{\alpha}\|^2$
since $f(\mathbf{X}_s) = \sum_{s'} \alpha_{s'} k(\mathbf{X}_{s'}, \mathbf{X}_s)
= [\mathbf{K}\boldsymbol{\alpha}]_s$.
For the RKHS norm,
\[
\|f\|_{\mathcal{H}}^2
= \left\langle \sum_s \alpha_s\, k(\mathbf{X}_s, \cdot),\;
         \sum_{s'} \alpha_{s'}\, k(\mathbf{X}_{s'}, \cdot) \right\rangle_{\mathcal{H}}
= \sum_{s,s'} \alpha_s\, \alpha_{s'}\, k(\mathbf{X}_s, \mathbf{X}_{s'})
= \boldsymbol{\alpha}^\top \mathbf{K}\, \boldsymbol{\alpha}.
\]
By condition~\eqref{eq:penalty_eval}, each $C_j(f)$ is a function of
$\hat{\mathbf{f}} = \mathbf{K}\boldsymbol{\alpha}$.
Combining yields~\eqref{eq:finite_dim}.
\end{proof}

%----------------------------------------------------------------------
\subsection{Proof of Proposition~\ref{prop:foc} (First-Order Conditions)}
%----------------------------------------------------------------------

\begin{proof}
Write the objective~\eqref{eq:finite_dim} with quadratic penalties as
\begin{multline}
\label{eq:obj_expanded}
L(\boldsymbol{\alpha})
= \frac{1}{n}(\mathbf{R} - \mathbf{K}\boldsymbol{\alpha})^\top
  (\mathbf{R} - \mathbf{K}\boldsymbol{\alpha})
+ \lambda_{\mathrm{stat}}\,\boldsymbol{\alpha}^\top\mathbf{K}\boldsymbol{\alpha} \\
+ \sum_{j \in \mathcal{J}^A} \mu_j
  (\mathbf{A}_j\mathbf{K}\boldsymbol{\alpha} - \mathbf{b}_j)^\top
  (\mathbf{A}_j\mathbf{K}\boldsymbol{\alpha} - \mathbf{b}_j)
+ \sum_{j \in \mathcal{J}^{B,C}} \frac{\mu_j}{n}
  (\mathbf{K}\boldsymbol{\alpha} - \mathbf{v}_j)^\top
  (\mathbf{K}\boldsymbol{\alpha} - \mathbf{v}_j),
\end{multline}
where $\mathbf{v}_j = \mathbf{g}_j$ for Type~B and $\mathbf{v}_j = \mathbf{h}_j$ for Type~C.
Taking the gradient with respect to $\boldsymbol{\alpha}$:
\begin{align}
\frac{\partial L}{\partial \boldsymbol{\alpha}}
&= -\frac{2}{n}\,\mathbf{K}(\mathbf{R} - \mathbf{K}\boldsymbol{\alpha})
+ 2\lambda_{\mathrm{stat}}\,\mathbf{K}\boldsymbol{\alpha}
+ 2\sum_{j \in \mathcal{J}^A} \mu_j\,\mathbf{K}\mathbf{A}_j^\top
  (\mathbf{A}_j\mathbf{K}\boldsymbol{\alpha} - \mathbf{b}_j)
\notag \\
&\quad + 2\sum_{j \in \mathcal{J}^{B,C}} \frac{\mu_j}{n}\,\mathbf{K}
  (\mathbf{K}\boldsymbol{\alpha} - \mathbf{v}_j).
\label{eq:gradient}
\end{align}
Setting~\eqref{eq:gradient} to zero and left-multiplying by $\mathbf{K}^{-1}$
(which exists when $\mathbf{K} \succ 0$):
\[
\frac{1}{n}\,\mathbf{K}\boldsymbol{\alpha}
+ \lambda_{\mathrm{stat}}\,\boldsymbol{\alpha}
+ \sum_{j \in \mathcal{J}^A} \mu_j\,\mathbf{A}_j^\top\mathbf{A}_j
  \mathbf{K}\boldsymbol{\alpha}
+ \sum_{j \in \mathcal{J}^{B,C}} \frac{\mu_j}{n}\,\mathbf{K}\boldsymbol{\alpha}
= \frac{1}{n}\,\mathbf{R}
+ \sum_{j \in \mathcal{J}^A} \mu_j\,\mathbf{A}_j^\top\mathbf{b}_j
+ \sum_{j \in \mathcal{J}^{B,C}} \frac{\mu_j}{n}\,\mathbf{v}_j.
\]
Collecting terms using the definitions of $\boldsymbol{\Omega}$
and $\boldsymbol{\omega}$ from~\eqref{eq:Omega}--\eqref{eq:omega_vec}
gives~\eqref{eq:foc}.
Strict positive definiteness of $\mathbf{K}$ and $\lambda_{\mathrm{stat}} > 0$
ensure the coefficient matrix in~\eqref{eq:foc} is invertible, yielding
uniqueness.
\end{proof}

%----------------------------------------------------------------------
\subsection{Proof of Proposition~\ref{prop:nesting} (Nesting)}
%----------------------------------------------------------------------

\begin{proof}
Claims~(i) and~(ii) follow directly from the definition.
For~(i), setting $\mu_j = 0$ eliminates the structural penalties.
For~(ii), taking $\mu_j \to \infty$ enforces $C_j(f) = 0$ as a constraint;
formally, the sequence of penalized minimizers converges to the
constrained minimizer provided $C_j$ is lower semicontinuous and
$\{f \in \mathcal{H} : C_j(f) = 0\}$ is nonempty.
Setting $\lambda_{\mathrm{stat}} = 0$ removes agnostic regularization.

For~(iii), the linear kernel restricts $\mathcal{H}$ to linear functions
$f(\mathbf{x}) = \boldsymbol{\beta}^\top \mathbf{x}$, with
$\|f\|_{\mathcal{H}}^2 = \|\boldsymbol{\beta}\|^2$.
Substituting a Type~C penalty with target-PCA factor loadings and
simplifying gives~\eqref{eq:ipca_nest}.

For~(iv), the correspondence is notational:
the $\ell_1$ norm replaces the $\ell_2$ RKHS norm,
and the structural penalty becomes a weighted $\ell_1$ term on coefficients
rather than a penalty on function evaluations.
The estimator~\eqref{eq:llm_lasso_nest} coincides with the adaptive
Lasso when $\gamma_{j,k} = |\hat{\beta}_k^{\mathrm{init}}|^{-1}$
and with text-informed penalization when $\gamma_{j,k}$ reflects
LLM-derived variable relevance scores.
\end{proof}

%----------------------------------------------------------------------
\subsection{Proof of Proposition~\ref{prop:misspecification} (Misspecification Robustness)}
%----------------------------------------------------------------------

\begin{proof}
We show that any $\mu_j > 0$ introduces an asymptotic bias, so the
oracle sets $\mu_j^* = 0$.

\emph{Step 1: Asymptotic bias from misspecification.}
When $C_j(f^*) > 0$, the set
$\mathcal{F}_j(\mu_j) = \arg\min_{f \in \mathcal{H}}
\frac{1}{n}\sum_s (R_s - f(\mathbf{X}_s))^2
+ \lambda_{\mathrm{stat}}\|f\|_{\mathcal{H}}^2 + \mu_j C_j(f)$
converges (as $n \to \infty$ and $\lambda_{\mathrm{stat}} \to 0$ at an
appropriate rate) to the population minimizer
\[
f_j^*(\mu_j) = \arg\min_{f} \mathbb{E}[(R - f(\mathbf{X}))^2] + \mu_j C_j(f).
\]
By~\ref{cond:consistency}, $f_j^*(0) = f^*$.
For $\mu_j > 0$, the term $\mu_j C_j(f)$ tilts the solution toward
$\{f : C_j(f) \text{ is small}\}$.
Since $C_j(f^*) > 0$, this tilt is binding: $f_j^*(\mu_j) \neq f^*$
for all $\mu_j > 0$, and the prediction risk satisfies
\[
\mathbb{E}\bigl[(R - f_j^*(\mu_j)(\mathbf{X}))^2\bigr]
> \mathbb{E}\bigl[(R - f^*(\mathbf{X}))^2\bigr]
= \mathbb{E}[\varepsilon^2]
\qquad \text{for all } \mu_j > 0.
\]

\emph{Step 2: Risk gap is bounded away from zero.}
By continuity of $C_j$ (condition~\ref{cond:bounded}) and strict
convexity of the population prediction risk,
for any $\delta > 0$ there exists $c(\delta) > 0$ such that
\[
\mathrm{Risk}(\mu_j) - \mathrm{Risk}(0)
\;\geq\; c(\delta) > 0
\qquad \text{for all } \mu_j \geq \delta,
\]
where $\mathrm{Risk}(\mu_j) = \mathbb{E}[(R - f_j^*(\mu_j)(\mathbf{X}))^2]$.

\emph{Step 3: Oracle efficiency implies $\hat{\mu}_j \to 0$.}
By~\ref{cond:cv_oracle}, $\mathrm{CV}(\hat{\boldsymbol{\lambda}})$
tracks the oracle $\inf_{\boldsymbol{\lambda}} \mathrm{CV}(\boldsymbol{\lambda})$
up to $o_p(1)$.  Since the oracle selects $\mu_j^* = 0$
(the unique risk minimizer along the $\mu_j$ coordinate by Step~2),
and the risk gap is bounded below by $c(\delta)$ for $\mu_j \geq \delta$,
we have $\Pr(\hat{\mu}_j \geq \delta) \to 0$ for every $\delta > 0$.
\end{proof}

%----------------------------------------------------------------------
\subsection{Proof of Proposition~\ref{prop:test_mu} (Testing the Marginal Value of Economic Models)}
%----------------------------------------------------------------------

\begin{proof}
\emph{Under $H_0$.}
When $\mu_j^* = 0$, the oracle solution does not use model $j$.
Constraining $\mu_j = 0$ therefore imposes no binding restriction on the
population objective.
The difference
$\mathrm{CV}(\hat{\boldsymbol{\lambda}}_{-j}) -
\mathrm{CV}(\hat{\boldsymbol{\lambda}})$
arises solely from finite-sample noise in the cross-validated selection
of the remaining hyperparameters.
By condition~\ref{cond:test_rate} and the oracle efficiency
in~\ref{cond:cv_oracle}, this difference is
$O_p(n^{-1} r_n)$, so $T_j = n \cdot O_p(n^{-1} r_n) = O_p(r_n) = o_p(1)$.

\emph{Under $H_1$.}
When $\mu_j^* > 0$, constraining $\mu_j = 0$ removes a binding
restriction.
By condition~\ref{cond:test_smooth},
the population risk increases by at least
$c > 0$ (a positive constant depending on $\mu_j^*$ and $C_j$).
By uniform convergence of the cross-validated risk to the population risk,
$\mathrm{CV}(\hat{\boldsymbol{\lambda}}_{-j}) -
\mathrm{CV}(\hat{\boldsymbol{\lambda}}) \geq c/2 > 0$
with probability approaching one.
Hence $T_j = n(c/2 + o_p(1)) \to \infty$.
\end{proof}
